% !TeX program = xelatex
% Самодостаточный исходный файл. Сборка: xelatex mathematical_review.tex
% Для разрешения ссылок запустите XeLaTeX дважды.
\documentclass[12pt,a4paper]{article}
\usepackage[margin=24mm,top=23mm,bottom=24mm,headheight=16pt]{geometry}
\usepackage{fontspec}
\defaultfontfeatures{Ligatures=TeX}
\setmainfont{Liberation Serif}
\setsansfont{Liberation Sans}
\setmonofont{DejaVu Sans Mono}[Scale=MatchLowercase]
\usepackage{polyglossia}
\setdefaultlanguage{russian}
\setotherlanguage{english}
\newfontfamily\cyrillicfont{Liberation Serif}
\newfontfamily\cyrillicfontsf{Liberation Sans}
\newfontfamily\cyrillicfonttt{DejaVu Sans Mono}[Scale=MatchLowercase]
\usepackage{amsmath,amssymb}
\usepackage{array,booktabs,tabularx,longtable}
\usepackage{enumitem}
\usepackage{microtype}
\usepackage{fancyhdr}
\usepackage{titlesec}
\usepackage{xurl}
\usepackage[unicode,hidelinks]{hyperref}
\hypersetup{
  pdftitle={Математическая рецензия на диссертацию Абдуллы Дуаа},
  pdfsubject={Корректность результатов, научная новизна и замечания к диссертации},
  pdfauthor={},
  pdfkeywords={комбинаторика слов, символическая динамика, рецензия, PhD}
}
\pagestyle{fancy}
\fancyhf{}
\fancyhead[L]{\small\itshape Математическая рецензия}
\fancyhead[R]{\small Абдулла Дуаа}
\fancyfoot[C]{\thepage}
\renewcommand{\headrulewidth}{0.3pt}
\setlength{\parindent}{1.25em}
\setlength{\parskip}{3pt}
\linespread{1.06}
\setlist[itemize]{leftmargin=1.6em,itemsep=4pt,topsep=5pt}
\setlist[enumerate]{leftmargin=1.7em,itemsep=5pt,topsep=5pt}
\titleformat{\section}{\large\bfseries}{\thesection.}{0.55em}{}
\titleformat{\subsection}{\normalsize\bfseries}{\thesubsection.}{0.55em}{}
\titlespacing*{\section}{0pt}{16pt plus 3pt minus 2pt}{7pt}
\titlespacing*{\subsection}{0pt}{11pt plus 2pt minus 1pt}{5pt}
\setcounter{secnumdepth}{2}
\emergencystretch=2em
\clubpenalty=10000
\widowpenalty=10000
\displaywidowpenalty=10000
\allowdisplaybreaks[1]
\newcommand{\irr}{\operatorname{irr}}
\newcommand{\ind}{\operatorname{ind}}
\newcommand{\dens}{\operatorname{dens}}
\newcommand{\sref}[1]{\textnormal{[Д, с.~#1]}}

\begin{document}
\thispagestyle{plain}
\begin{center}
  {\LARGE\bfseries Математическая рецензия\par}
  \vspace{8pt}
  {\large на диссертацию Абдуллы Дуаа\par}
  \vspace{5pt}
  {\large «Комбинаторика слов и символическая динамика:\par
  теоретические основы, плотность слова Фибоначчи,\par
  практические применения и топологические индексы»\par}
  \vspace{8pt}
  {Экспертная оценка с позиции комиссии по присуждению степени PhD\par}
  \vspace{4pt}
  {17 ноября 2026 года\par}
\end{center}
\vspace{4pt}

\noindent\textbf{Итоговая рекомендация:} вернуть работу на существенную
доработку перед присуждением степени PhD. В диссертации имеется
содержательное математическое ядро, прежде всего результаты о скрученных
повторах и деревьях факторов. Однако представленная редакция содержит
явные математические ошибки, недостаточно точное обоснование новизны
и несоразмерно большой объём вспомогательного материала.

Рецензия относится к представленной рукописи 2026 года
\cite{dissertation}, PDF объёмом 149 страниц, и сопровождающему её
исходному тексту. В ссылках на диссертацию используется
\emph{печатная нумерация страниц}, а не порядковый номер страницы PDF.
Обозначение \sref{49}, например, означает страницу 49 рукописи.
Номера теорем и формул также относятся к диссертации, если явно
не указано иное.

\section{Общая математическая оценка}

Работа объединяет несколько направлений: распределение букв и факторов
в слове Фибоначчи, палиндромные префиксы, повторы в обобщённых словах
Туэ--Морса, графовые индексы и педагогический эксперимент. Первые четыре
направления допускают общую исследовательскую постановку. Педагогическая
глава связана с ними преимущественно тематически.

Положительной чертой является явное признание автором классического
происхождения многих результатов. В частности, значение палиндромной
плотности слова Фибоначчи и формула для характеристических штурмовых слов
правильно приписаны Фишлеру. Эти результаты действительно содержатся
в его работе 2006 года \cite{fischler}.

Для оценки диссертации необходимо различать корректное самостоятельное
изложение известной теории, новые элементарные следствия и новые
результаты, определяющие исследовательский вклад.

\begin{longtable}{@{}>{\raggedright\arraybackslash}p{0.46\textwidth}
                      >{\raggedright\arraybackslash}p{0.50\textwidth}@{}}
\toprule
\textbf{Результат} & \textbf{Оценка}\\
\midrule
\endfirsthead
\toprule
\textbf{Результат} & \textbf{Оценка}\\
\midrule
\endhead
\bottomrule
\endfoot
Полная факторная сложность конкатенации факториалов, \S\,2.2
& Доказательство выглядит корректным; арифметическая сложность
следует непосредственно.\\[5pt]
Точный подсчёт букв и сбалансированность, \S\,3.4
& Корректное изложение классических свойств.\\[5pt]
Погрешность плотностей обрамлённых слов, \S\,3.5.4
& Корректное, но небольшое уточнение.\\[5pt]
Богатство тернарного слова, \S\,3.6.3
& Корректный результат, непосредственно связанный
с вставкой разделителя.\\[5pt]
Классификация скрученных повторов, теорема 3.8.5
& Наиболее содержательный претендент
на основной новый результат.\\[5pt]
Динамические свойства, теорема 3.8.7
& В основном стандартные следствия примитивности подстановки.\\[5pt]
Индексы деревьев факторов, теорема 5.1.3
& Содержательный точный расчёт; доказательство выглядит согласованным.\\
\end{longtable}

\section{Оценка основных исследовательских результатов}

\subsection{Классификация скрученных повторов}

Теорему 3.8.5 \sref{94--95} следует считать наиболее сильной частью
работы. Для
\[
  W^{(b)}[n]=s_b(n)\pmod N
\]
в непериодическом случае $N\nmid b-1$ утверждается существование сильного
$(k,\sigma^j)$-повтора в точности при
\[
  k\le \kappa(N,b,j):=
  \max\left\{2,\left\lfloor\frac{2b}{j'}\right\rfloor\right\},
  \qquad j'\in\{1,\ldots,N\},\qquad j'\equiv j\pmod N.
\]
Здесь $s_b(n)$ --- сумма цифр числа $n$ в системе счисления
с основанием $b$, а $\sigma$ --- циклическая перестановка алфавита.

Доказательство имеет ясный механизм: условие повторения переводится
в постоянство разности сумм цифр; множители $b$ в длине блока
устраняются; затем анализ переносов ограничивает оставшуюся длину.
Отдельно строятся повторы, достигающие заявленной границы.
Существенного разрыва в этой цепочке рассуждений при проведённой
проверке не обнаружено.

Однако литературное сопоставление необходимо усилить. Критический
показатель обычных повторов в тех же обобщённых словах Туэ--Морса
вычислен Blondin-Massé, Brlek, Glen и Labbé в 2007 году
\cite{blondin}. При $j=0$ рассматриваемая классификация согласуется
с целочисленной частью их результата. Поэтому новый вклад следует
формулировать прежде всего для \emph{нетождественного скручивания}
и подробно объяснить его отношение к известным методам.

Имеется локальная ошибка в последней строке доказательства:
из $j'\ge b$ следует
\[
  \left\lfloor\frac{2b}{j'}\right\rfloor\le2,
\]
а не $\le1$. При $j'=b$ значение равно $2$.
Итоговая формула от этого не меняется.

\subsection{Индексы деревьев факторов}

Теорема 5.1.3 \sref{128--129} представляет второй убедительный
исследовательский фрагмент. Существенно, что автор рассматривает точно
определённое дерево всех факторов длины не более $n$.

Ключевой шаг доказательства --- описание смежности правых специальных
факторов. Если характеристическое слово начинается с $x^dy$, то соседние
специальные вершины соединены лишь в начальной цепочке. При $n\ge d+2$
это позволяет получить числа рёбер соответствующих типов:

\begin{center}
\begin{tabular}{@{}cc@{}}
\toprule
\textbf{Степени концов} & \textbf{Число рёбер}\\
\midrule
$1,3$ & $2$\\
$1,2$ & $n-1$\\
$2,3$ & $3n-3-2d$\\
$3,3$ & $d-1$\\
\bottomrule
\end{tabular}
\end{center}

Отсюда непосредственно получаются
\[
  \irr(T_n)=4n-2d,\qquad
  \irr_2(T_n)=18n-10d-2.
\]
Подсчёты для индексов Загреба
\[
  M_1(T_n)=2n^2+8n-4,\qquad
  M_2(T_n)=2n^2+10n-11+d
\]
также согласуются с указанной структурой.

Результат имеет ясный смысл: рассматриваемые индексы отражают длину
начальной серии характеристического слова. Его глубину следует оценивать
соразмерно: после структурного наблюдения вычисления элементарны,
а последующие неполные частные углового коэффициента в эти формулы
не входят.

\subsection{Дополнительные конечные проверки}

Формулы теоремы 5.1.3 дополнительно проверены на 20 независимо
построенных конечных деревьях для нескольких штурмовых угловых
коэффициентов; расхождений не обнаружено.

Конечный поиск скрученных повторов при $2\le b\le8$, $2\le N\le9$,
$N\nmid b-1$, длинах блоков до $64$ в префиксах длины $20\,000$
также не дал нарушений заявленной верхней границы. Эти вычисления
дополняют проверку доказательств, но не заменяют её и не доказывают
достижимость каждой границы в рассмотренных конечных префиксах.

\section{Существенные математические замечания}

\subsection{Распределение логарифмов факториалов}

Перед леммой 2.2.1 \sref{43} говорится, что последовательность
$\log(n!)$ не является равномерно распределённой по модулю $1$.
Это утверждение неверно: для каждого фиксированного основания $b>1$
последовательность
\[
  \{\log_b(n!)\}
\]
равномерно распределена на $[0,1)$. Для десятичного основания этот
результат, в частности, доказан в статье Диакониса о распределении
ведущих цифр \cite{diaconis}.

Общий случай получается из формулы Стирлинга и оценки Ван дер Корпута.
Для главной гладкой части
\[
  f(x)=\frac{(x+\tfrac12)\log x-x}{\log b}
\]
имеем
\[
  f''(x)=\frac{x^{-1}-(2x^2)^{-1}}{\log b}\asymp\frac1x.
\]
Это даёт необходимое затухание нормированных экспоненциальных сумм
в критерии Вейля.

Собственная лемма автора о бесконечном числе попаданий в каждый
интервал остаётся полезным элементарным доказательством достаточного
для дальнейшего рассуждения свойства. Ошибочную характеристику
известной последовательности необходимо удалить.

\subsection{Противоречивое определение обобщённой последовательности}

В определении 3.1.2 \sref{49} одновременно заданы
\[
  F_0^p(t,k)=\cdots=F_{k-2}^p(t,k)=0,
  \qquad F_{k-p-1}^p(t,k)=1.
\]
При $k=3$, $p=1$ первое условие даёт $F_1^1(t,3)=0$,
а второе --- $F_1^1(t,3)=1$.

Таким образом, определение в приведённой форме не задаёт математического
объекта даже для простых допустимых параметров. Требуется восстановить
правильные ограничения на параметры, начальные условия и диапазон
рекурсии, а затем согласовать их с леммой 3.1.1.

\subsection{Несогласованная индексация в формуле (3.4)}

Две суммы представлены как выражения одного $F_n$ \sref{49}:
\[
  F_n=\sum_{k=0}^{\lfloor(n-1)/2\rfloor}\binom{n-1-k}{k},
  \qquad
  F_n=\sum_{k=0}^{\lfloor n/2\rfloor}\binom{n-k}{k}.
\]
При принятой далее нормировке $F_0=0$, $F_1=1$ вторая сумма равна
$F_{n+1}$. Уже при $n=3$ суммы равны соответственно $2$ и $3$.
Индексацию необходимо исправить.

\subsection{Некорректная формулировка теоремы 3.1.1}

Начальное предложение теоремы \sref{50--51} грамматически
и математически незавершено; величины $\ind$ и $\ind^{*}$
не определены. Формула
\[
  \ind^{*}(u)=2\limsup_{n\to\infty}[\,a_n,a_{n-1},\ldots,a_1\,]
\]
требует исправления.

Если имеется в виду асимптотический критический показатель,
стандартная формула имеет вид
\[
  E^{*}(u)=2+\limsup_{n\to\infty}[\,a_n,a_{n-1},\ldots,a_1\,].
\]
Для слова Фибоначчи это даёт $2+\varphi$, тогда как напечатанное
выражение даёт $2\varphi$. Формула с суммой приведена, например,
у Adamczewski--Allouche \cite[разд.~4.3, замечание~11]{adamczewski}.

Необходимы сверка с первоисточником и точное определение всех
используемых индексов. При другом смысле обозначения $\ind^{*}$
этот смысл должен быть явно установлен до формулировки теоремы.

\subsection{Статус задачи о факторной сложности}

В наблюдении 3.8.3 \sref{93} и заключении закономерность
\[
  p(k+1)-p(k)\in\{N(N-1),N^2\}
\]
оставлена как вычислительное наблюдение, требующее доказательства.

Между тем Tromp и Shallit ещё в 1995 году определили полную функцию
факторной сложности именно последовательности $s_2(n)\bmod N$
\cite{tromp}. Следовательно, данное наблюдение необходимо сопоставить
с известной формулой и вывести из неё либо точно объяснить, какая
дополнительная задача остаётся нерешённой.

\subsection{Гипотеза о максимальности индекса звезды}

Автор выдвигает гипотезу \sref{126}:
\[
  \irr_2(T)<\irr_2(S_n)\qquad\text{при }T\not\cong S_n
\]
и сообщает о проверке при $n\le12$. Здесь
\[
  \irr_2(T)=\sum_{uv\in E(T)}|d_u^2-d_v^2|,
\]
а $S_n$ --- звезда на $n$ вершинах.

Доказательство получается простой оценкой. У дерева, отличного
от звезды, максимальная степень $\Delta\le n-2$, а степени концов
любого ребра не меньше $1$. Поэтому
\[
  |d_u^2-d_v^2|\le\Delta^2-1.
\]
Поскольку дерево имеет $n-1$ ребро,
\[
\begin{aligned}
  \irr_2(T)
  &\le(n-1)(\Delta^2-1)\\
  &\le(n-1)\bigl((n-2)^2-1\bigr)\\
  &<(n-1)\bigl((n-1)^2-1\bigr)
   =\irr_2(S_n).
\end{aligned}
\]
Таким образом, гипотезу следует заменить предложением
с доказательством и исправить соответствующий пункт заключения.

\section{Масштаб и обоснование научной новизны}

Значительная часть главы 3 --- точные частоты букв, сбалансированность,
описание через последовательности Битти --- представляет стандартную
теорию штурмовых слов. Самостоятельные доказательства ценны для
связности изложения, но их нельзя учитывать как новые исследовательские
достижения.

Тернарное слово
\[
  \mathbf y=\iota_0(\rho(\mathbf f)),\qquad \iota_0(c)=0c,
\]
является прозрачным кодированием слова Фибоначчи. Лемма о соответствии
палиндромов и формула
\[
  |\mathrm P(0x_10x_2\cdots0x_m0)|=1+2|\mathrm P(x)|
\]
выглядят корректными. Значимость результата определяется общим
свойством операции вставки разделителя; конкретная тернарная
конструкция служит его примером. Новизну следует обсуждать
в контексте морфизмов, сохраняющих богатство.

Минимальность, строгая эргодичность и нулевая энтропия
в теореме 3.8.7 в основном следуют из стандартной теории примитивных
подстановок. Их доказательство для выбранного семейства полезно,
но само по себе не составляет самостоятельного нового динамического
результата.

Предложенная в \S\,3.7 метрика является обычной суммируемой
взвешенной метрикой произведения конечных пространств.
Для самостоятельного исследовательского результата необходимо показать,
какую новую задачу она позволяет решить.

Теорема о разрезании полимино корректна в выбранной модели полных
линий сетки без перемещения частей. При таком определении необходимость
каждой линии, содержащей внутреннее ребро, фактически встроена
в постановку. Это небольшой вспомогательный результат.

Точное выражение погрешности для обрамлённых слов Фибоначчи
представляет корректное уточнение. Его следует оценивать как
элементарное следствие точных количеств букв и формулы Бине.
Перечисление большого числа таких следствий не должно затруднять
выделение основных результатов диссертации.

\section{Педагогическая глава и воспроизводимость}

Автор обоснованно признаёт исходное неравенство групп и отсутствие
статистически значимого скорректированного эффекта. Такая осторожность
выводов заслуживает положительной оценки.

Однако использование целых классов требует обсуждения внутриклассной
зависимости и единицы статистического анализа; приведённого анализа
на уровне учеников для этого недостаточно.

Для диссертации по математике эта глава не увеличивает вес
теоретического вклада. Её разумно существенно сократить либо
вынести в приложение.

В рукописи сообщается о внешнем комплекте программ и журнале
из 121 проверки. Указанный комплект при подготовке рецензии получить
не удалось; поэтому заявленные автором 121 проверка независимо
не воспроизведены. Это ограничение доступа само по себе не является
свидетельством отсутствия программ. Для повторного рецензирования
следует приложить к представляемой версии самодостаточный архив
программ, входных данных и журналов вычислений.

\clearpage
\section{Необходимые изменения и итоговое заключение}

Для положительного решения комиссии необходимо:
\begin{enumerate}
  \item Исправить перечисленные математические ошибки и повторно
  проверить предварительные определения и формулы.
  \item Сосредоточить положения, выносимые на защиту, вокруг
  действительно новых результатов, прежде всего теорем 3.8.5
  и 5.1.3.
  \item Провести подробное сопоставление с литературой о критических
  показателях, факторной сложности и морфизмах, сохраняющих богатство.
  \item Заменить элементарно разрешимые открытые вопросы доказанными
  утверждениями и актуализировать заключение.
  \item Приложить программы и журналы вычислений в форме,
  позволяющей воспроизвести заявленные проверки.
\end{enumerate}

\noindent\textbf{Итоговое заключение.}
Работа содержит потенциально достаточное для дальнейшего рассмотрения
исследовательское ядро, однако её математическая выверенность
и обоснование объёма нового вклада пока не соответствуют требованиям
к окончательной диссертации PhD. Рекомендуется существенная переработка
с последующим повторным рецензированием.
\emph{В представленной редакции присуждение степени не рекомендуется.}

\begin{thebibliography}{9}
\small
\bibitem[Д]{dissertation}
Абдулла Дуаа.
\emph{Комбинаторика слов и символическая динамика:
теоретические основы, плотность слова Фибоначчи,
практические применения и топологические индексы}.
Диссертация. Московский физико-технический институт.
Москва, 2026. Представленная рукопись, 149 страниц PDF.

\bibitem{fischler}
Fischler S.
Palindromic prefixes and episturmian words.
\emph{Journal of Combinatorial Theory, Series A}.
2006. Vol.~113. P.~1281--1304.
\url{https://www.imo.universite-paris-saclay.fr/~stephane.fischler/palindromes.pdf}.

\bibitem{blondin}
Blondin-Massé A., Brlek S., Glen A., Labbé S.
On the critical exponent of generalized Thue--Morse words.
\emph{Discrete Mathematics \& Theoretical Computer Science}.
2007. Vol.~9, no.~1. P.~293--304.
\url{https://doi.org/10.46298/dmtcs.397}.

\bibitem{diaconis}
Diaconis P.
The distribution of leading digits and uniform distribution mod~1.
\emph{The Annals of Probability}.
1977. Vol.~5, no.~1. P.~72--81.
\url{https://doi.org/10.1214/aop/1176995891}.

\bibitem{adamczewski}
Adamczewski B., Allouche J.-P.
Reversals and palindromes in continued fractions.
\emph{Theoretical Computer Science}.
2007. Vol.~380. P.~220--237.
\url{https://doi.org/10.1016/j.tcs.2007.03.017}.

\bibitem{tromp}
Tromp J., Shallit J.
Subword complexity of a generalized Thue--Morse word.
\emph{Information Processing Letters}.
1995. Vol.~54, no.~6. P.~313--316.
\url{https://doi.org/10.1016/0020-0190(95)00074-M}.
\end{thebibliography}
\end{document}
